% Insert in the simulation section: association-module comparison setup.
% Existing main Algorithm 1 and the training procedure are unchanged.
% Requires the existing change color and the accompanying BibTeX entries.
\subsubsection{\textcolor{change}{Association-Module Comparison}}
\textcolor{change}{
To assess the performance--complexity tradeoff of the distance-priority
association, we compare it with two literature-inspired alternatives,
Matching and Swap. Only the association module is replaced; both
alternatives retain the proposed resource allocation and energy safety layer.
}

\textcolor{change}{
\emph{Matching:} Following the many-to-one request-and-rejection mechanism
in \cite{sun2024twotimescale}, each active user ranks the UAVs within its
coverage by an estimated latency-reduction utility. For preference
construction only, let $\tilde f_{i,k}=F_{\mathrm{UAV}}/N_{\max}$ and
$\tilde b_{i,k}=W/N_{\max}$. Using the original offloading proposal
$\hat x_i(t)$, the pairwise utility is
}
\begin{equation}
\textcolor{change}{
 u_{i,k}(t)=1-
 \frac{\max\{\widehat T_i^{\mathrm{loc}}(t),
                   \widetilde T_{i,k}^{\mathrm{edge}}(t)\}}
      {T_i^{\mathrm{base}}(t)},
}
\label{eq:association_preference}
\end{equation}
\textcolor{change}{
where $\widehat T_i^{\mathrm{loc}}(t)$ is the local-processing latency
under $\hat x_i(t)$, and $\widetilde T_{i,k}^{\mathrm{edge}}(t)$ is the
corresponding edge-path latency under the nominal resources.
Users apply to UAVs in descending order of positive utility. Each UAV
retains up to $N_{\max}$ applicants with the highest utility and rejects
the others, which subsequently apply to their next preferred UAVs.
The preferences remain fixed during matching, and unmatched users
execute locally. The pricing utility in the reference is replaced by
\eqref{eq:association_preference} to match our objective.
}

\textcolor{change}{
\emph{Swap:} Inspired by the two-phase matching strategy in
\cite{wang2025matching}, a distance-based association first provides
reference UAV loads. A rate-based matching is then constructed using
bandwidth estimates $W/\max\{1,n_k^{\mathrm{dist}}(t)\}$, where
$n_k^{\mathrm{dist}}(t)$ is the number of users associated with UAV $k$
in the distance-based initialization. Starting from this matching,
feasible user exchanges are evaluated, including transfers to spare
capacity and local execution through dummy options. Each candidate
undergoes resource allocation and energy correction. An exchange is
accepted only if it decreases
}
\begin{equation}
\textcolor{change}{
 J(\mathbf A;t)=\sum_{i:\xi_i(t)=1}
 \frac{\max\{T_i^{\mathrm{loc}}(\mathbf A;t),
             T_i^{\mathrm{edge}}(\mathbf A;t)\}}
      {T_i^{\mathrm{base}}(t)},
}
\label{eq:association_evaluation}
\end{equation}
\textcolor{change}{
where $\mathbf A=[a_{i,k}]$ denotes the association matrix; all active
users, including unmatched users executing locally, are included.
The first feasible improvement exceeding $10^{-8}$ is accepted,
followed by a new neighborhood scan. The search stops when no
improving neighbor is found or the budget of $20$ rounds or $1000$
candidate evaluations is reached. Budget-limited termination is
reported separately. These alternatives adapt the matching mechanisms
of the references to our model and do not reproduce their complete
optimization frameworks.
}

\textcolor{change}{
All comparisons use a frozen DARS policy. At each evaluated state,
the candidate UAV positions and original offloading proposals are
shared across the three methods; association masking and energy
correction are recomputed separately. Only the original DARS decisions
advance the environment. We report the mean SLRR, paired SLRR gains,
and mean and 95th-percentile module runtimes. Runtime includes
preference construction, matching or exchange search, and all
associated resource-allocation and safety checks. Thus, the comparison
isolates instantaneous association effects under the frozen policy's
state distribution, rather than establishing joint global optimality
or the long-term performance of separately trained alternatives.
}
